\section{Rajoitukset}
Tutkimukseen liittyy useita rajoituksia, jotka on huomioitava tuloksia tulkittaessa.

Ensinnäkin tutkimus kohdistuu yhteen järjestelmään, Snowflaken Cortex Analystiin, jonka taustalla toimii Anthropicin Claude Sonnet 4 -malli. Tulokset kuvaavat siten yhtä tiettyä toteutusta eivätkä ole suoraan yleistettävissä muihin kielimalleihin, SQL-avustajiin tai luonnollisen kielen analytiikkatyökaluihin.

Toiseksi tutkimus perustuu yhteen sovellusalueeseen eli puhelinkeskusdataan. Tämän tietokannan rakenne, sanasto ja analyysitarpeet voivat poiketa merkittävästi muiden toimialojen vastaavista ympäristöistä. On mahdollista, että järjestelmä suoriutuisi eri tavoin esimerkiksi talousdatan, lokidatan tai asiakastietojen yhteydessä.

Kolmanneksi tutkimuksessa tarkastellaan mallin tuottamia SQL-kyselyitä eikä kyselyiden palauttamia varsinaisia vastauksia. Tämä rajaus on perusteltu aineiston luottamuksellisuuden vuoksi, mutta samalla se tarkoittaa, että tutkimus ei mittaa suoraan liiketoiminnallisen lopputuloksen oikeellisuutta. Kysely voi olla rakenteeltaan puutteellinen tai vaihtoehtoisesti syntaktisesti erilainen mutta käytännössä toimiva, eikä tätä eroa voida aina arvioida täysin ilman varsinaista tulosaineistoa.

Neljänneksi semanttisen näkymän laatu vaikuttaa suoraan siihen, miten hyvin Cortex Analyst kykenee muodostamaan kyselyitä. Jos näkymässä on epäselvyyksiä, puutteita tai monitulkintaisia kuvauksia, osa havaituista ongelmista voi johtua itse määrittelystä eikä pelkästään kielimallin rajoituksista. Tämän vuoksi tutkimus arvioi samalla myös sitä, kuinka hyvin kyseinen semanttinen näkymä tukee luonnollisen kielen ja tietokantarakenteen välistä tulkintaa.

Viidenneksi kielimallien toimintaa luonnehtii tietty epädeterministisyys, minkä vuoksi tulokset voivat vaihdella eri suorituskerroilla tai järjestelmäversioiden välillä. Lisäksi kaupalliset kielimallipohjaiset järjestelmät voivat muuttua ajan myötä ilman, että kaikki muutokset ovat käyttäjälle näkyviä. Tämän vuoksi tutkimuksen havainnot kuvaavat järjestelmän toimintaa tiettynä ajankohtana ja tietyssä konfiguraatiossa eivätkä välttämättä pysyvää ominaisuutta \cite{chen2024chatgpt}.

Kuudenneksi tuotettujen SQL-kyselyiden oikeellisuuden arvioi yksi arvioija. Tämä heikentää luokittelun reliabiliteettia ja lisää mahdollisuutta siihen, että osa luokitteluista sisältää arvioijan tulkintaan liittyvää subjektiivisuutta.

\section{Johtopäätökset}

Tämän tutkielman tavoitteena oli arvioida, kuinka hyvin Snowflaken Cortex Analyst soveltuu luonnollisen kielen tietokantakyselyihin puhelinkeskusdataa sisältävässä tietokantaympäristössä. Tarkastelun kohteena olivat erityisesti järjestelmän kyky tuottaa oikeellisia SQL-kyselyitä, sen konsistenssi tilanteissa, joissa sama kysymys esitetään useita kertoja, sekä sen herkkyys sellaisille syötteen muotoilun muutoksille, joiden ei pitäisi muuttaa kysymyksen merkitystä.

Tulokset osoittavat, että Cortex Analyst kykenee monissa tapauksissa tuottamaan oikeellisia ja keskenään johdonmukaisia SQL-kyselyitä. Erityisesti yksinkertaisemmat kysymykset sekä sellaiset kysymykset, joihin oli määritelty valmiita metriikoita semanttisessa näkymässä, tuottivat hyvin vakaita tuloksia. Tämä viittaa siihen, että kielimallipohjainen luonnollisen kielen rajapinta voi toimia käyttökelpoisena välineenä tietokannan hyödyntämisessä silloin, kun kysymykset ovat rakenteeltaan selkeitä ja semanttinen näkymä tukee niitä riittävän hyvin.

Tutkimus osoitti kuitenkin myös merkittäviä rajoituksia. Järjestelmän suorituskyky heikkeni erityisesti silloin, kun kysymykseen sisältyi implisiittisiä taustaoletuksia tai kun oikean SQL-kyselyn muodostaminen edellytti tarkkaa suodatusta, duplikaattien hallintaa tai monitulkintaisen sanamuodon tulkintaa. Nämä havainnot ovat linjassa aiemman Text-to-SQL-tutkimuksen kanssa, jossa vastaavien monimutkaisten ehtojen ja taustatietojen yhdistämisen on todettu olevan kielimalleille huomattava haaste \cite{li2024bird,yu2018spider}. Lisäksi havaittiin, että semanttisesti lähes saman sisältöiset kysymykset saattoivat johtaa erilaisiin SQL-kyselyihin ja joissakin tapauksissa myös erilaisiin virheisiin. Tämä ilmentää kielimalleille tyypillistä syöteherkkyyttä (\textit{prompt sensitivity}), jossa pienetkin sanamuodon vaihtelut voivat muuttaa tuotettua lopputulosta merkittävästi \cite{chatterjee2024posix}, mikä heikentää luonnollisen kielen rajapinnan ennustettavuutta käyttäjän näkökulmasta.

Keskeinen havainto on, että järjestelmän luotettavuus ei riipu yksinomaan taustalla olevasta kielimallista, vaan myös siitä, miten hyvin semanttinen näkymä on suunniteltu. Tutkimuksessa havaitut virheet liittyivät usein tilanteisiin, joissa tietokannan kenttien merkitysten erottelu ei ollut mallin näkökulmasta riittävän selkeä. Tämän perusteella voidaan päätellä, että semanttinen näkymä ei ole pelkästään tekninen apurakenne, vaan keskeinen osa koko järjestelmän toimintaa. Sen huolellinen suunnittelu, yksiselitteinen nimeäminen ja iteratiivinen kehittäminen voivat parantaa luonnollisen kielen kyselyjärjestelmän luotettavuutta merkittävästi.

Tutkimuksen perusteella konsistenssia ei myöskään voida pitää yksinään riittävänä laadun mittarina. Järjestelmä saattoi tuottaa saman tuloksen toistuvasti myös silloin, kun tulos oli virheellinen tai kun SQL-kyselyä ei tuotettu lainkaan. Toisaalta joissakin tapauksissa samaan kysymykseen saatiin useita erilaisia SQL-kyselyitä, joista osa oli oikeellisia ja osa virheellisiä. Tämä osoittaa, että luonnollisen kielen tietokantarajapintojen arvioinnissa on tarkasteltava samanaikaisesti sekä oikeellisuutta että vaihtelua.

Kokonaisuutena tutkielma tukee alan yleistä näkemystä, jonka mukaan luonnollisen kielen käyttö tietokantojen rajapintana on lupaava mutta vielä osin epävarma lähestymistapa \cite{liu2024survey,deng2022recent}. Cortex Analyst voi toimia tehokkaana työkaluna erityisesti rajatuissa käyttötapauksissa, joissa tietokannan rakenne tunnetaan hyvin ja semanttinen näkymä on laadittu huolellisesti. Sen sijaan tilanteissa, joissa kysymykset sisältävät implisiittisiä oletuksia, monitulkintaisuutta tai useita päättelyvaiheita, järjestelmän tuottamiin kyselyihin on syytä suhtautua varauksella.

Tutkielman tulokset viittaavat siihen, että käytännön kehitystyössä huomio kannattaa kohdistaa paitsi mallin suorituskykyyn myös siihen, miten järjestelmälle tarjottu semanttinen konteksti on rakennettu. Tulevassa tutkimuksessa olisi hyödyllistä vertailla useita erilaisia semanttisia näkymiä, testata järjestelmää muissa sovellusympäristöissä sekä tarkastella tuotettujen SQL-kyselyiden lisäksi myös niiden palauttamien vastausten oikeellisuutta. Lisäksi useamman arvioijan käyttö voisi parantaa luokittelun reliabiliteettia ja vahvistaa tulosten luotettavuutta.

Tämän tutkimuksen perusteella voidaan päätellä, että kielimallipohjaiset luonnollisen kielen tietokantarajapinnat eivät vielä täysin korvaa asiantuntijan laatimia kyselyitä, mutta ne voivat jo nyt toimia hyödyllisenä tukena analytiikkatyössä. Niiden käytännön arvo riippuu kuitenkin ratkaisevasti siitä, kuinka hyvin järjestelmän toimintaympäristö, semanttinen malli ja käyttötapa on sovitettu yhteen.
